% Active manuscript insertion: reduced lab-scale survival closure.
% The exact finite-chain sections remain the microscopic reference layer.

\section{Reduced laboratory-scale survival closure}
\label{sec:reduced_survival_closure}

The exact finite-chain probability projection is retained as a microscopic reference, but it is not autonomous in the initial one-point spacing marginal because microscopic ordering and correlation information is lost under projection. Laboratory-frequency audits further show that conservative normal dynamics approaches a quasistatic cycle on the atomic mechanical time scale. Therefore high-cycle accumulation is formulated as survivor first passage rather than as an imposed permanent drift of the normalized spacing density.

Let $P_0(\lambda)$ denote the reference-phase structural/prestress spacing density, $q(t)=\sigma(t)/E$, and $q_{\rm ref}$ the reference-phase reduced traction. Each starting spacing is embedded by
\begin{equation}
q_r(\lambda_0)=\phi'(\lambda_0)-q_{\rm ref},
\end{equation}
so that its intact quasistatic trajectory satisfies
\begin{equation}
\phi'[\Lambda(\lambda_0,t)]
=\phi'(\lambda_0)+q(t)-q_{\rm ref},
\qquad
\phi''(\Lambda)>0.
\end{equation}
Consequently
\begin{equation}
\dot\Lambda=\frac{\dot q(t)}{\phi''(\Lambda)}.
\end{equation}
The corresponding nonabsorbing structural density obeys
\begin{equation}
\partial_tP+\partial_\lambda\left[
\frac{\dot q(t)}{\phi''(\lambda)}P
\right]=0.
\end{equation}
Thus the reversible mechanical part is closed directly from $P_0$ and the applied stress history under the declared quasistatic/prestress reduction.

The operational instability boundary remains
\begin{equation}
\lambda_c=\left(\frac{m+1}{n+1}\right)^{1/(m-n)},
\qquad \phi''(\lambda_c)=0.
\end{equation}
For a stable spacing $\lambda$, define the effective-potential climb
\begin{equation}
\Delta\psi_c(\lambda)
=\left[\phi(\lambda_c)-\phi'(\lambda)\lambda_c\right]
-\left[\phi(\lambda)-\phi'(\lambda)\lambda\right].
\end{equation}
If $A_c$ is the characteristic cohesive area of one local event, the barrier is
\begin{equation}
\Delta G_c=EA_ca_0\Delta\psi_c.
\end{equation}
Under the explicitly declared rare-event and fast intrawell thermal re-equilibration assumptions, the positive-flux transition-state rate is
\begin{equation}
k_c(\lambda,T;A_c)
=\frac{\sqrt{\phi''(\lambda)}}{2\pi t_0}
\exp\left[-\frac{EA_ca_0\Delta\psi_c(\lambda)}{k_BT}\right].
\label{eq:kc_reduced}
\end{equation}
No fitted diffusion coefficient, damping constant, lifetime distribution, or scalar damage variable is introduced in Eq.~\eqref{eq:kc_reduced}. The additional physical assumptions are finite temperature, rapid re-equilibration inside the intact well, and rare transition-state crossing.

The intact survivor subdensity therefore obeys the closed reduced equation
\begin{equation}
\boxed{
\partial_tP_b
+\partial_\lambda\left[
\frac{\dot\sigma(t)/E}{\phi''(\lambda)}P_b
\right]
=-k_c(\lambda,T;A_c)P_b,
}
\label{eq:final_survivor_pde}
\end{equation}
with $P_b(\lambda,t_0)=P_0(\lambda)$. Hence
\begin{equation}
S(t)=\int P_b(\lambda,t)d\lambda,
\qquad
F_{\rm ci}(t)=1-S(t).
\end{equation}
This gives the reduced mapping
\begin{equation}
\boxed{
P_0+\sigma(0:t)\longrightarrow P_b(\lambda,t),\ S(t),\ F_{\rm ci}(t).
}
\end{equation}

For periodic loading, each starting spacing has an integrated one-cycle hazard
\begin{equation}
\mathcal H_c(\lambda_0)=\int_0^{T_f}k_c[\Lambda(\lambda_0,t),T;A_c]dt,
\end{equation}
and therefore
\begin{equation}
\boxed{
S_N=\int P_0(\lambda_0)e^{-N\mathcal H_c(\lambda_0)}d\lambda_0.
}
\label{eq:cycle_survival}
\end{equation}
Permanent cycle-to-cycle drift of the normalized structural PDF is therefore not required. The cumulative history is stored in survivor mass, while the surviving normalized population may change only by selection.

The characteristic area $A_c$ is deliberately not fitted in the present derivation. A sensitivity audit at the current historical Al bridge, $300\,$K, and a $100\pm100\,$MPa, $20\,$Hz sinusoid gives one-cycle escape probabilities of approximately $4.66\times10^{-3}$, $2.30\times10^{-6}$, and $1.16\times10^{-9}$ for illustrative area ratios $A_c/A_0=40,50,60$, respectively. These values demonstrate exponential sensitivity only; no area ratio is adopted. In particular, a local per-cycle survival extremely close to unity is the expected high-cycle regime rather than a contradiction.

In the strict quasistatic phase-controlled limit, $\mathcal H_c\propto1/f$. The reduced hypothesis therefore predicts approximately frequency-independent life in physical time and a cycle count approximately proportional to loading frequency until the quasistatic or local-equilibrium assumptions fail. This, together with the temperature, mean-stress, stress-amplitude, and $P_0$ dependence of Eq.~\eqref{eq:final_survivor_pde}, makes the closure experimentally falsifiable.

The reduced equation is mathematically closed but is not asserted to be a complete dislocation theory of real single-crystal aluminum fatigue. Independent determination of $A_c$ and structural $P_0$, validation of the operational dividing surface $\lambda_c$, and any transition-state recrossing correction remain experimental or higher-fidelity tasks. Specimen-scale correlation area/volume and local-to-specimen survival scaling are deliberately deferred to a later calibration stage.